\begin{chapterpage}{Device Drivers}
    \begin{chapteritems}
        \item Introduction
        \item Driver API
    \end{chapteritems}
\end{chapterpage}

\section{Introduction}

X16 Open is based on loadable device drivers that do most of the work in
the system.

Device drivers needed to get the system up and running are embedded into
the kernel ROM bank, and are automatically registered and initialized when
the system is reset.

All drivers can be replaced at runtime, making it possible to
customize the system.


\section{Device Driver Dispatch Function}

All device drivers are required to implement a dispatch function that is placed at the
start of the driver's place in memory.

The dispatch function takes a "sub-function" number in the X registers. Most sub-functions
need more arguments that are passed in the A, Y, and arg1-10 registers as 
described below.

The following sub-functions must be supported:

\begin{longtblr}{colspec = {X[1,l] X[2,l] X[4,l]},
width = \linewidth,
rowhead = 1,
rowfoot = 0}
\# & Name    & Description \\
\hline
0   & on\_register      & Called when the device driver is registered \\
1   & on\_unregister    & Called when the device driver is unregistered \\
2   & open              & Opens a file \\
3   & read\_byte        & Reads one byte from a file \\
4   & read\_block       & Reads one block of bytes from a file \\
5   & write\_byte       & Writes one byte to a file \\
6   & write\_block      & Writes a block of bytes to a file \\
7   & seek              & Moves the file pointer \\
8   & position          & Returns the file pointer \\
9   & close             & Closes a file \\
10  & ioctl             & Driver control command \\
11  & on\_irq           & Called on every IRQ \\
\end{longtblr}


\section{Function: on\_register}

\textbf{Driver function \#:} 0 \\
\textbf{Input:} None \\
\textbf{Output:} A = Flags \\
\textbf{Errors:} EIO

This function is called by the kernel when registering a device driver. It can,
for instance, be used to initialize the driver.

The device driver flags (bits 0-5) are returned in the A register. The possible flags are:

\begin{itemize}
    \item Bit 0: Set if reading is supported.
    \item Bit 1: Set if writing is supported.
    \item Bit 2: Set if appending is supported.
    \item Bit 3: Set if seeking is supported.
    \item Bit 4: Set if a block device, unset if a char device.
    \item Bits 5-6: Reserved, set by the kernel, not the driver.
    \item Bit 7: Set if kernel must call the driver's on\_irq function on every IRQ.
\end{itemize}

EIO is a driver-specific error.


\section{Driver Function: on\_unregister}

\textbf{Driver function \#:} 1 \\
\textbf{Input:} None \\
\textbf{Output:} None  \\
\textbf{Errors:} EIO

This function is called by the kernel when unregistering a device driver. It can,
for instance, be used to clean up and free driver resources.

EIO is a driver-specific error.


\section{Function: open}

\textbf{Function \#:} 2 \\
\textbf{Input:} A = Flags, arg1-2 = Path address, arg3 = Path bank \\
\textbf{Output:} Y = Context ID, A = Flags \\
\textbf{Errors:} ENOTFOUND, ENOTFILE, EINVAL, EIO

Opens a file or directory at the specified path.

You must select exactly one of the following flags with
the A register:

\begin{itemize}
    \item Bit 0: Set if open for reading
    \item Bit 1: Set if open for writing
    \item Bit 2: Set if open for appending
\end{itemize}

The driver's internal context ID is returned in the Y register if the call is successful. 

The A register holds the following flags:

\begin{itemize}
    \item Bits 0-2: Input preserved.
    \item Bit 3: Set if a file, unset if a directory.
    \item Bit 7: Set if opened successfully.
\end{itemize}

ENOTFOUND is returned if the file is not found.

ENOTFILE is returned if attempting to write or append
to a directory.

EINVAL is returned if the operation requested by the input flags is 
not supported by the device, or if not exactly one of the flags was set.

EIO is a driver-specific error.


\section{Function: read\_byte}

\textbf{Function \#:} 3 \\
\textbf{Input:} Y = Context ID \\
\textbf{Output:} A = Byte value \\
\textbf{Errors:} ENOTOPEN, EEOF, EIO

Reads one byte from a file.

The context ID is used internally by the device driver to keep
track of multiple open files.

ENOTOPEN is returned if the context ID is not active.

EEOF is returned if no value can be returned, as the file
pointer is at the end of the file. EEOF is not returned when
successfully reading the final byte; it is returned on
the following read attempt.

EIO is a driver-specific error.


\section{Function: read\_block}

\textbf{Function \#:} 4 \\
\textbf{Input:} Y = Context ID, arg1-2 = Requested number of bytes, arg3-4 = Buffer address, arg5 = Buffer bank \\
\textbf{Output:} arg6-7 = Number of bytes read \\
\textbf{Errors:} ENOTOPEN, EEOF, EIO

Reads one block of bytes from a file. The requested size of the block is set in arg1-2.
The block is stored at the memory address and bank set in arg 3-5.

The size of a block is driver-specific. It is not guaranteed that you receive the requested 
number of bytes. The number of bytes actually read is returned in arg6-7.

The context ID is used internally by the device driver to keep
track of multiple open files.

ENOTOPEN is returned if the context ID is not active.

EEOF is returned if no value can be returned, as the file
pointer is at the end of the file. EEOF is not returned when
successfully reading the final byte; it is returned on
the following read attempt.

EIO is a driver-specific error.


\section{Function: write\_byte}

\textbf{Function \#:} 5 \\
\textbf{Input:} Y = Context ID, A = Byte value \\
\textbf{Output:} None \\
\textbf{Errors:} ENOTOPEN, ENOTFILE, EIO

Writes one byte to a file.

The context ID is used internally by the device driver to keep
track of multiple open files.

ENOTOPEN is returned if the context ID is not active.

ENOTFILE is returned if attempting to write or append
to a directory.

EIO is a driver-specific error.


\section{Function: write\_block}

\textbf{Function \#:} 6 \\
\textbf{Input:} Y = Context ID, arg1-2 = Buffer size, arg3-4 = Buffer address, arg5 = Buffer bank \\
\textbf{Output:} arg6-7 = Number of bytes written \\
\textbf{Errors:} ENOTOPEN, ENOTFILE, EIO

Writes one block of bytes to a file.

The size of a block is driver-specific. It is not guaranteed that all requested
bytes are written to the file. The number of bytes actually written is returned in arg6-7.

The context ID is used internally by the device driver to keep
track of multiple open files.

ENOTOPEN is returned if the context ID is not active.

ENOTFILE is returned if attempting to write or append
to a directory.

EIO is a driver-specific error.


\section{Function: seek}

\textbf{Function \#:} 7 \\
\textbf{Input:} Y = Context ID, arg1-4 = New position \\
\textbf{Output:} arg1-4 = Current position \\
\textbf{Errors:} EINVAL, ENOTOPEN, EIO

Moves the file pointer to the requested position.

The file position after this call is returned in arg1-4.

Seeking is only supported on files opened for reading. Attempting to 
seek on a file opened for writing or appending returns 
the EINVAL error.

EINVAL is also returned if the device driver doesn't support
seeking.

The ENOTOPEN error is returned if the specified file context ID is
not active.

EIO is a driver-specific error.


\section{Function: position}

\textbf{Function \#:} 8 \\
\textbf{Input:} Y = Context ID \\
\textbf{Output:} arg1-4 = File position \\
\textbf{Errors:} EINVAL, ENOTOPEN, EINVAL, EIO

Returns the current position in the specified file in arg1-4.

ENOTOPEN is returned if the context ID is not active.

EINVAL can be returned if the driver doesn't support
reading the file positions.

EIO is a driver-specific error.


\section{Function: close}

\textbf{Function \#:} 9 \\
\textbf{Input:} Y = Context ID \\
\textbf{Output:} None \\
\textbf{Errors:} ENOTOPEN, EIO

This function closes a file.

ENOTOPEN is returned if the context ID is not active.

EIO is a driver-specific error.


\section{Function: ioctl}

\textbf{Function \#:} 10 \\
\textbf{Input:} Y = Context ID, A = Device type, arg1 = Command, arg2 = Input length, arg3-4 = Input data or input buffer address, arg5 = Input buffer bank,
arg6 = Output max length, arg7-8 = Output buffer address, arg9 = Output buffer bank \\
\textbf{Output:} A = Output length, arg7-8 = Output data or output buffer address, arg9 = Output buffer bank \\
\textbf{Errors:} EIO

This function executes a driver-specific control command.

The file context ID is passed in the Y register. The driver may ignore the context ID if it is not relevant
to the control command.

The device type is passed in the A register. A driver for a specific device type is expected to
handle certain control commands.

The control command is specified in arg1. The two most significant bits are used to flag if the input 
(arg3-4) and output (arg7-8) are direct values or buffer pointers. Direct values can be at most two bytes long.

\begin{itemize}
    \item Bit 7: Set if output is a direct value, clear if output is a buffer pointer.
    \item Bit 6: Set if input is a direct value, clear if input is a buffer pointer.
\end{itemize}

The input is specified in arg2-5:

\begin{itemize}
    \item arg2: Input length.
    \item arg3-4: A direct value (arg1 bit 6 set) set by the driver or input buffer address set by the caller (arg bit 6 clear).
    \item arg5: Input buffer bank set by the caller (arg1 bit 6 clear)
\end{itemize}

The output is specified in arg6-9 and the A register:

\begin{itemize}
    \item arg6: Output max length set by the caller.
    \item arg7-8: A direct value (arg 1 bit 7 set) set by the driver or output buffer address set by the caller (arg1 bit 7 clear).
    \item arg9: Output buffer bank set by the caller (arg1 bit 7 clear).
    \item A: Output length set by the driver.
\end{itemize}

EIO is a driver-specific error.


\section{Function: on\_irq}

\textbf{Function \#:} 11 \\
\textbf{Input:} None \\
\textbf{Output:} None \\
\textbf{Errors:} EIO

This function is only intended to be called by the kernel's
IRQ handler.

It can be used by device drivers that need to be updated
once every interrupt, for instance terminals.

EIO is a driver-specific error.
